\section{\texorpdfstring{$\lambda$}{Lambda}-specifications}

Throughout, $S$ is a set, $(E, \mathcal{E})$ a measurable space and $\nu$ a measure on $E$.

\begin{definition}[Product probability measure]
    \label{def:product-probability-measure}
    \lean{MeasureTheory.Measure.infinitePi}
    \leanok{}

    For probability spaces $(\Omega_i,\MCB_i,P_i)$, $i \in I$, the product
    $P:=\bigotimes_{i\in I}P_i$ is a probability measure on $\prod_{i\in I}\Omega_i$.
\end{definition}

\begin{definition}[Independent Specification with Single Spin Distribution (ISSSD)]
    \label{def:isssd}
    \uses{def:juxtaposition}
    \lean{Specification.isssdFun}
    % \leanok

    For $\nu \in \mathfrak{P}(E, \mathcal{E})$, the \textbf{Independent Specification with Single
    Spin Distribution $\nu$} is
    \begin{align}
        \ISSSD:\mathfrak{P}(E, \mathcal{E})&\to\Finset(S)\to\mathcal{E}^S\times E^S\to\overline{\mathbb{R}_{\geq0}}\\
        \nu&\mapsto\Lambda\mapsto(A\mid\omega)\mapsto\left(\nu^\Lambda\left(\juxtBy{\omega}^{-1}(A)\right)\right)
    \end{align}
    where $\nu^\Lambda$ is the product measure on $E^\Lambda$.
\end{definition}

\begin{lemma}[Independence of ISSSDs]
    \label{lem:isssd-indep}
    \uses{def:isssd, def:indep-spec}
    \lean{Specification.isssdFun_comp_isssdFun}
    \leanok{}

    $\ISSSD(\nu)_{\Lambda_1} \circ_k \ISSSD(\nu)_{\Lambda_2} = \ISSSD(\nu)_{\Lambda_1 \union \Lambda_2}$
    for all $\Lambda_1, \Lambda_2 \in \Finset(S)$.
\end{lemma}
\begin{proof}
    % \leanok

    Both sides agree on square cylinders, a $\pi$-system generating $\mathcal{E}^S$.
\end{proof}


\begin{definition}[ISSSDs are specifications]
    \label{def:isssd-spec}
    \uses{lem:isssd-indep, def:spec}
    \lean{Specification.isssd}
    \leanok

    $\ISSSD(\nu)$ is a specification.
\end{definition}

\begin{lemma}[ISSSDs are proper specifications]
    \label{lem:isssd-proper-spec}
    \uses{def:isssd, def:markov-spec}
    \lean{Specification.IsProper.isssd}
    \leanok

    $\ISSSD(\nu)$ is a proper specification.
\end{lemma}
\begin{proof}
    \uses{def:isssd-spec}
    % \leanok

    An event of $\mathcal{F}_{S\setminus\Lambda}$ is determined by $\omega$ outside $\Lambda$, so
    it factors out of $\nu^\Lambda(\juxtBy{\omega}^{-1}(\cdot))$.
\end{proof}

\subsection*{Free measures}

\begin{definition}[Free measure of a specification]
    \label{def:free-measure}
    \uses{def:spec}
    \lean{Specification.HasFreeMeasure}
    \leanok

    A measure $\mu_0$ on $(\Cfg, \mathcal{F})$ is a \textbf{free measure} of a specification
    $\gamma$ if $\gamma_\Lambda(A \mid \eta) = \mu_0(A)$ for every $\Lambda \in \Fin$, every
    $\eta \in \Cfg$ and every $A \in \mathcal{F}_\Lambda$.
\end{definition}

\begin{theorem}[Free measures determine the Gibbs measures]
    \label{thm:free-measure-gibbs}
    \uses{def:free-measure, def:gibbs-meas}
    \lean{Specification.eq_smul_of_hasFreeMeasure, Specification.eq_of_hasFreeMeasure,
      Specification.isGibbsMeasure_iff_indep_of_hasFreeMeasure}
    \leanok

    Let $\gamma$ have free measure $\mu_0 \in \Prob(\Cfg, \mathcal{F})$. Every finite
    $\mu \in \Gibbs{\gamma}$ equals $\mu(\Cfg)\,\mu_0$; in particular $\mu_0$ is the only
    probability measure that can lie in $\Gibbs{\gamma}$. Moreover $\mu_0 \in \Gibbs{\gamma}$ if
    and only if $\mathcal{F}_\Lambda$ and $\mathcal{F}_{S \setminus \Lambda}$ are independent
    under $\mu_0$ for every $\Lambda \in \Fin$.
\end{theorem}
\begin{proof}
    \leanok

    For $A \in \mathcal{F}_\Lambda$ the DLR equation gives
    $\mu(A) = \int \gamma_\Lambda(A \mid \eta)\,\mu(\mathrm{d}\eta) = \mu(\Cfg)\,\mu_0(A)$, and
    the local events form a $\pi$-system generating $\mathcal{F}$. For the second assertion, let
    $A \in \mathcal{F}_\Lambda$ and $B \in \mathcal{F}_{S \setminus \Lambda}$; properness gives
    $\gamma_\Lambda(A \inter B \mid \eta) = 1_B(\eta)\,\mu_0(A)$, so
    $\mu_0\gamma_\Lambda(A \inter B) = \mu_0(A)\,\mu_0(B)$, which equals $\mu_0(A \inter B)$
    exactly under the stated independence. Conversely a square cylinder splits as
    $A \inter B$ with $A$ over the sites in $\Lambda$ and $B$ over those outside, and square
    cylinders generate.
\end{proof}

\begin{lemma}[Uniqueness of a Gibbs measure specified by an ISSSD]
    \label{lem:isssd-gibbs-meas-uniqueness}
    \uses{def:gibbs-meas, def:isssd-spec, thm:free-measure-gibbs}
    \lean{Specification.isGibbsMeasure_isssd_iff,
      Specification.isGibbsMeasure_isssdFamily_iff}
    \leanok

    A probability measure $\mu$ is a Gibbs measure for $\ISSSD(\nu)$ if and only if
    $\mu = \nu^S$; likewise, for a family $(\nu_i)_{i \in S}$ of probability measures, $\mu$ is a
    Gibbs measure for the inhomogeneous $\ISSSD$ of $(\nu_i)$ if and only if
    $\mu = \bigotimes_{i \in S} \nu_i$.
\end{lemma}
\begin{proof}
    \leanok

    $\ISSSD(\nu)$ has free measure $\nu^S$, under which $\mathcal{F}_\Lambda$ and
    $\mathcal{F}_{S \setminus \Lambda}$ are independent; apply
    Theorem \ref{thm:free-measure-gibbs}.
\end{proof}

\begin{lemma}[Existence of a Gibbs measure specified by an ISSSD]
    \label{lem:isssd-gibbs-meas-existence}
    \uses{def:gibbs-meas, def:product-probability-measure, def:isssd-spec}
    \lean{Specification.isGibbsMeasure_isssd_infinitePi}
    \leanok{}

    $\nu^S$ is a Gibbs measure for $\ISSSD(\nu)$.
\end{lemma}
\begin{proof}
    % \leanok

    On square cylinders both sides of the DLR equation factor over the coordinates.
\end{proof}

\begin{definition}[Modifier]
    \label{def:modif}
    \uses{def:spec}
    \lean{Specification.IsModifier, Specification.modification}
    \leanok{}

    A {\bf modifier of $\gamma$} is a family
    \[\rho : \Finset(S) \to \Omega \to [0, \infty[\]
    of measurable functions such that the kernels
    $(\rho\gamma)_\Lambda(\cdot \mid \eta) := \rho_\Lambda \cdot \gamma_\Lambda(\cdot \mid \eta)$
    form a specification $\rho\gamma$.
\end{definition}

\begin{lemma}[Modifier of a modifier]
    \label{lem:modif-modif}
    \uses{def:modif}
    \lean{Specification.modification_modification}
    \leanok{}

    If $\rho_1$ is a modifier of $\gamma$ and $\rho_2$ a modifier of $\rho_1\gamma$, then
    $\rho_1\rho_2$ is a modifier of $\gamma$ and $\rho_2(\rho_1\gamma) = (\rho_1\rho_2)\gamma$.
\end{lemma}
\begin{proof}
    % \uses{}
    \leanok

    Both sides have density $\rho_{1,\Lambda}\rho_{2,\Lambda}$ with respect to
    $\gamma_\Lambda(\cdot\mid\eta)$.
\end{proof}

\begin{lemma}[A modifier of a proper specification is proper]
    \label{lem:modif-proper}
    \uses{def:modif, def:markov-spec}
    \lean{Specification.isProper_modification}
    \leanok{}

    If $\rho$ is a modifier of a specification $\gamma$, then $\rho\gamma$ is proper.
\end{lemma}
\begin{proof}
    \uses{lem:proper-ker-lintegral}
    \leanok

    For $\Lambda \in \Finset(S)$, $A \in \mathcal E^S$, $B \in \mathcal{F}_{S\setminus\Lambda}$
    and $\eta \in \Cfg$,
    \[(\rho\gamma)_\Lambda(A \inter B \mid \eta)
      = \int_{A \inter B} \rho_\Lambda \,\mathrm{d}\gamma_\Lambda(\cdot\mid\eta)
      = 1_B(\eta) \int_{A} \rho_\Lambda \,\mathrm{d}\gamma_\Lambda(\cdot\mid\eta)
      = 1_B(\eta)\,(\rho\gamma)_\Lambda(A \mid \eta)\]
    by Lemma \ref{lem:proper-ker-lintegral} with $f = 1_A\rho_\Lambda$, $g = 1_B$.
\end{proof}

\begin{lemma}[Null events of Gibbs measures, Remark 1.28.2]
    \label{lem:remark-128-2}
    \uses{def:modif, def:gibbs-meas}
    \lean{Specification.IsGibbsMeasure.lambdaSpecification_null_iff,
      Specification.IsGibbsMeasure.lambdaSpecification_null_iff_null}
    \leanok

    Let $\nu$ be a $\sigma$-finite non-zero measure on $E$ with $\lambda$-kernels
    $\lambda_\Lambda(\cdot \mid \eta) = \nu^\Lambda \otimes \delta_{\eta_{S \setminus \Lambda}}$,
    let $h$ be a premodifier (Definition \ref{def:premodif}) with $h_\Lambda > 0$ everywhere
    and $0 < \lambda_\Lambda h_\Lambda(\eta) < \infty$ for all $\Lambda \in \Fin$, $\eta \in \Cfg$,
    and let $\gamma = \rho\lambda$ with $\rho_\Lambda = h_\Lambda / \lambda_\Lambda h_\Lambda$.
    A finite non-zero $\mu \in \Gibbs{\gamma}$ annihilates an event $C \in \mathcal{F}_\Delta$,
    $\Delta \in \Fin$, exactly when $\lambda_\Delta(C \mid \eta) = 0$ for every $\eta$. In
    particular any two finite non-zero Gibbs measures for $\gamma$ have the same null events in
    $\bigcup_{\Delta \in \Fin} \mathcal{F}_\Delta$.
\end{lemma}
\begin{proof}
    \leanok

    The DLR equation at the volume $\Delta$ writes $\mu(C)$ as the $\mu$-average of
    $\gamma_\Delta(C \mid \cdot) = \int_C \rho_\Delta \, \mathrm{d}\lambda_\Delta(\cdot\mid\eta)$,
    and positivity of $\rho_\Delta$ makes the inner integral vanish exactly when
    $\lambda_\Delta(C \mid \eta) = 0$. On $\mathcal{F}_\Delta$ the value
    $\lambda_\Delta(C \mid \eta)$ does not depend on $\eta$.
\end{proof}

\begin{lemma}[Every specification is a modification of some ISSSD, Remark 1.28.5]
    \label{lem:exists-modif-isssd}
    \uses{def:modif, def:isssd}
    % \lean{}
    % \leanok

    If $E$ is countable, $\nu$ is the counting measure and $\gamma$ is any specification, then
    $$\rho_\Lambda(\eta) = \gamma_\Lambda(\{\sigma_\Lambda = \eta_\Lambda\} | \eta)$$
    is a modifier from $\ISSSD(\nu)$ to $\gamma$.
\end{lemma}
\begin{proof}
    % \uses{}
    % \leanok

    Fix $\Lambda \in \Finset(S)$ and $\eta \in \Cfg$. The measure $\ISSSD(\nu)_\Lambda(\cdot\mid\eta)$
    is counting measure on the countable set $\set{\zeta : \zeta_{S \setminus \Lambda} = \eta_{S\setminus\Lambda}}$,
    and for such $\zeta$ properness gives
    $\rho_\Lambda(\zeta) = \gamma_\Lambda(\{\sigma_\Lambda = \zeta_\Lambda\} \mid \eta)$. Hence
    for measurable $A$
    \[(\rho\ISSSD(\nu))_\Lambda(A\mid\eta)
      = \sum_{\zeta_{S\setminus\Lambda} = \eta_{S\setminus\Lambda}} 1_A(\zeta)\,
        \gamma_\Lambda(\{\sigma_\Lambda = \zeta_\Lambda\} \mid \eta)
      = \gamma_\Lambda(A\mid\eta),\]
    since $\gamma_\Lambda(\cdot\mid\eta)$ is carried by the same countable set.
\end{proof}

\begin{proposition}[Characterisation of modifiers, Proposition 1.30.1]
    \label{prop:modif-char}
    \uses{def:modif}
    \lean{Specification.isModifier_iff_ae_eq}
    \leanok{}

    Let $\gamma$ be a specification and $\rho$ a family of functions $\Cfg \to [0,\infty]$. TFAE
    \begin{enumerate}
        \item $\rho$ is a modifier of $\gamma$
        \item each $\rho_\Lambda$ is measurable, $\int \rho_\Lambda \dd\gamma_\Lambda(\cdot|\eta) = 1$
        for every $\Lambda$ and $\eta$, and for all
        $\Lambda_1 \subseteq \Lambda_2$ and all $\eta \in \Cfg$,
        \[\rho_{\Lambda_2} = \rho_{\Lambda_1}\cdot (\gamma_{\Lambda_1} \rho_{\Lambda_2}) \quad \gamma_{\Lambda_2}(\cdot|\eta)\text{-a.e.}\]
    \end{enumerate}
    The normalisation in (2) cannot be dropped: $\rho_\Lambda \equiv 0$ satisfies the displayed
    identity and is not a modifier.
\end{proposition}
\begin{proof}
    % \uses{}
    % \leanok

    Normalisation says that each $(\rho\gamma)_\Lambda(\cdot|\eta)$ is a probability measure.
    Properness and consistency of $\gamma$ give, for $\Lambda_1 \subseteq \Lambda_2$,
    $(\rho\gamma)_{\Lambda_1} \circ_k (\rho\gamma)_{\Lambda_2}
    = \big(\rho_{\Lambda_1}\cdot(\gamma_{\Lambda_1}\rho_{\Lambda_2})\big)\gamma_{\Lambda_2}$, so
    consistency of $\rho\gamma$ equates two densities with respect to the same $\sigma$-finite
    measure $\gamma_{\Lambda_2}(\cdot|\eta)$, which is the displayed a.e. identity.
\end{proof}

\begin{proposition}[Characterisation of modifiers of independent specifications, Proposition 1.30.2]
    \label{prop:modif-indep-char}
    \uses{def:modif, def:indep-spec}
    \lean{Specification.isModifier_iff_ae_comm}
    \leanok{}

    Let $\gamma$ be a specification with
    $\gamma_{\Lambda_1} \circ_k \gamma_{\Lambda_2} = \gamma_{\Lambda_1 \union \Lambda_2}$ whenever
    $\Lambda_1 \inter \Lambda_2 = \emptyset$ (in particular any independent specification,
    Definition \ref{def:indep-spec}), and $\rho$ a family of functions $\Cfg \to [0,\infty]$.
    TFAE
    \begin{enumerate}
        \item $\rho$ is a modifier of $\gamma$
        \item each $\rho_\Lambda$ is measurable and normalised as in
        Proposition~\ref{prop:modif-char}, and for all $\Lambda_1 \subseteq \Lambda_2$,
        $\eta \in \Cfg$ and $\gamma_{\Lambda_2 \setminus \Lambda_1}(\cdot|\eta)$-almost all
        $\eta_2 \in \Cfg$,
        \[\rho_{\Lambda_2}(\zeta_1)\rho_{\Lambda_1}(\zeta_2) = \rho_{\Lambda_2}(\zeta_2) \rho_{\Lambda_1}(\zeta_1)\]
        for $\gamma_{\Lambda_1}(\cdot|\eta_2) \times \gamma_{\Lambda_2}(\cdot|\eta_2)$-almost all $(\zeta_1, \zeta_2)$.
    \end{enumerate}
\end{proposition}
\begin{proof}
    % \uses{}
    % \leanok

    By Proposition \ref{prop:modif-char} it suffices to compare the a.e. identity there with the
    symmetry in (2). Disjoint consistency gives
    $\gamma_{\Lambda_2}(\cdot|\eta) = \int \gamma_{\Lambda_1}(\cdot|\eta_2)\,
    \gamma_{\Lambda_2\setminus\Lambda_1}(\mathrm{d}\eta_2|\eta)$, so the identity holds
    $\gamma_{\Lambda_2}(\cdot|\eta)$-a.e. iff for $\gamma_{\Lambda_2\setminus\Lambda_1}(\cdot|\eta)$-almost
    all $\eta_2$ it holds $\gamma_{\Lambda_1}(\cdot|\eta_2)$-a.e. For fixed $\eta_2$ and
    normalised $\rho_{\Lambda_1}$, integrating the symmetry in $\zeta_2$ against
    $\gamma_{\Lambda_1}(\cdot|\eta_2)$ recovers
    $\rho_{\Lambda_2} = \rho_{\Lambda_1}\cdot(\gamma_{\Lambda_1}\rho_{\Lambda_2})$ a.e., and
    conversely that identity substituted on both sides gives the symmetry.
\end{proof}

\begin{definition}[Premodifier, Definition 1.31]
    \label{def:premodif}
    % \uses{}
    \lean{Specification.IsPremodifier}
    \leanok

    A family of measurable functions $h_\Lambda : \Cfg \to [0, \infty[$ is a {\bf premodifier} if
    $$h_{\Lambda_2}(\zeta)h_{\Lambda_1}(\eta) = h_{\Lambda_1}(\zeta)h_{\Lambda_2}(\eta)$$
    for all $\Lambda_1 \subseteq \Lambda_2$ and all $\zeta, \eta \in \Cfg$ with
    $\zeta_{S \setminus \Lambda_1} = \eta_{S \setminus \Lambda_1}$.
\end{definition}

\begin{lemma}[Modifiers are premodifiers]
    \label{lem:modif-premodif}
    \uses{def:modif, def:premodif}
    % \lean{}
    % \leanok
    If $\rho$ is a modifier of $\ISSSD(\nu^S)$, then it is a premodifier if any of the following conditions hold:
    \begin{enumerate}
        \item $E$ is countable and $\nu$ is equivalent to the counting measure.
        \item $E$ is a second countable Borel space.
        \item $\nu$ is everywhere dense.
        \item For all $\Lambda_1 \subseteq \Lambda_2$ and all $\eta \in \Cfg$, $\zeta \mapsto \rho_{\Lambda_1}(\zeta \eta_{S \setminus \Lambda_1})$ is continuous on $E^{\Lambda_1}$.
    \end{enumerate}
\end{lemma}
\begin{proof}
    \uses{prop:modif-char}
    % \leanok

    \begin{enumerate}
        \item Proposition \ref{prop:modif-char}: under counting measure, an a.e. identity holds
        everywhere.
        \item Omitted.
        \item Omitted.
        \item Omitted.
    \end{enumerate}
\end{proof}

\begin{lemma}[Premodifiers give rise to modifiers, Remark 1.32]
    \label{lem:premodif-modif}
    \uses{def:modif, def:premodif}
    \lean{Specification.IsPremodifier.isModifier_premodifierNorm}
    \leanok
    If $h$ is a premodifier, $\nu$ a probability measure and
    $0 < \ISSSD(\nu)_\Lambda h_\Lambda(\eta) < \infty$ for all $\Lambda$ and $\eta$, then
    $$\rho_\Lambda := \frac{h_\Lambda}{\ISSSD(\nu)_\Lambda h_\Lambda}$$
    is a modifier of $\ISSSD(\nu)$.
\end{lemma}
\begin{proof}
    \uses{prop:modif-indep-char}

    $\ISSSD(\nu)_\Lambda h_\Lambda$ is $\mathcal{F}_{S\setminus\Lambda}$-measurable, so
    $\ISSSD(\nu)_\Lambda\rho_\Lambda = 1$ and $\rho$ is again a premodifier. The premodifier
    identity is the pointwise form of the symmetry in
    Proposition \ref{prop:modif-indep-char}(2), since
    $\ISSSD(\nu)_\Lambda(\cdot|\eta) \times \ISSSD(\nu)_\Lambda(\cdot|\eta)$ is carried by pairs
    agreeing with $\eta$ outside $\Lambda$.
\end{proof}
